\section{The excitation--interaction matrix}\label{sec:9.1}

\Cref{tab:9-atlas} has 24 rows, one for each kind of excitation that this book treats, and six columns, one for each class of dark-matter coupling. Of its 144 cells, 57 hold a channel, and each of those names the one section that computes its rate. Eight more are allowed at leading order but computed in no section. The remaining 79 are empty at leading order, for reasons of operator content, kinematics or order in the coupling that we work out below. A channel is a pair of row and column, so the table is also a check that no channel is computed twice: where two sections seem to treat the same process, they treat different cells, or different parts of one response.

\subsection{Excitation and interaction axes}\label{sec:9.1-axes}

The columns come from the reduction of \cref{sec:8.1}. Every slow interaction, written in the operator basis $\mathcal{O}_1,\dots,\mathcal{O}_{15}$, reduces to a few target operators per species $s$: the number density $n_s$, the spin density $\vect{s}_s$, the paramagnetic current $\vect{j}_s$, the spin current $\mathcal{J}^{\alpha\beta}_s$ and the kinetic-energy density $\mathcal{T}_s$ (\cref{eq:8.1-rows,tab:8-operators}). Three distinctions separate the columns. The first is scattering against absorption. A scattered particle deposits $\omega$ over a region of the $(\vect{q},\omega)$ plane bounded by the kinematics of \cref{sec:1.3}, while an absorbed boson deposits its mass at the single point $\vect{q} = \vect{p}_\phi = m_\phi\vect{v}$, $\omega = m_\phi$ (\cref{sec:8.1-kinematics,sec:8.2-absorption}). The second is the species: the electron rows and the nuclear rows reach different final states, as \cref{sec:8.6} showed. The third is the operator: a density cannot flip a spin, and a spin density carries no charge.

\emph{Electron density} holds every scattering coupling whose leading target operator is $n_e$, for a heavy or a light mediator: $\mathcal{O}_1$ with the coefficient $c_1^{(e)} = \tilde V$, and the density parts of $\mathcal{O}_5$, $\mathcal{O}_8$ and $\mathcal{O}_{11}$. \emph{Nuclei} holds the couplings to $n_\kappa$, with the coherent weight $Zc_1^{(p)} + (A - Z)c_1^{(n)}$ of \cref{eq:8.1-nuclear}. \emph{Spin and current} holds the couplings to $\vect{s}_s$, $\vect{j}_s$ and $\mathcal{J}^{\alpha\beta}_s$, of electrons or of nuclei: $\mathcal{O}_4$ and the other operators of \cref{tab:8-operators} marked as flipping spins, and the current parts of $\mathcal{O}_2$, $\mathcal{O}_5$ and $\mathcal{O}_8$. In the lattice rows its cells are those of the nuclear spin. \emph{Vector absorption} is the absorption of a dark photon, whose row $-\vect{j}_e + \sum_\kappa Z_\kappa\vect{j}_\kappa$ is the current of all charges, or of another vector boson. \emph{Axion-like and scalar absorption} uses $\vect{s}_e$ and the trace $\mathcal{J}^{\alpha\alpha}_e$ for the axion-like particle (\cref{eq:8.1-alp}) and $n_e - \mathcal{T}_e/m_e$ for the scalar (\cref{eq:8.1-scalar}), or their nucleon analogues.

The last column keeps the target rows and changes the kinematics. It collects inelastic dark matter with a mass splitting $\Delta_\chi$ (\cref{eq:8.1-omega-inel}) and its exothermic version, which deposits more than its kinetic energy (\cref{sec:8.2-exothermic}). It also collects dark matter boosted to relativistic speed, and fermionic absorption, in which the dark matter turns into a light fermion and deposits energy along the line $\omega = \mchi - q$ (\cref{eq:8.1-omega-ferm}). A coupling to the electron density thus belongs to the first column when the dark matter scatters elastically and to the last when it scatters inelastically. The last column, however, holds a cell only where the change of kinematics changes the physics (\cref{sec:9.1-empty}).

The rows follow the solid-state sectors of \crefrange{chap:3}{chap:7}, in the order in which the channel chapters take them up. Electronic excitations at clamped nuclei come first, then lattice excitations, then the mixed processes in which a phonon and an electronic excitation are created together, and last the magnetic, superconducting and other targets. \Cref{tab:9-atlas} cites their defining sections. One row is new. \emph{Phonon-assisted electron recoil} is the creation of an electron--hole pair together with a phonon, which supplies crystal momentum that the dark matter cannot. It is the scattering counterpart of the phonon-assisted absorption of light in \cref{sec:6.5}, and \cref{sec:13.2}, which derives its second-order amplitude, makes it precise. In silicon it is the only way that dark matter scattering on electrons can create a pair at masses just below the zero-phonon threshold, because the gap is indirect (\cref{sec:3.3}). Rows that none of those chapters introduces carry a one-line description.

\begin{table}[tbp]
  \centering
  \footnotesize
  \setlength{\tabcolsep}{2.5pt}
  \renewcommand{\arraystretch}{1.05}
  \caption{The channel atlas. Rows: excitations of a solid, with their defining section or a short description. Sector: the final-state sector of \cref{eq:8.6-partition} (\cref{sec:8.6}). ad: electron-following; 0: zero-phonon electronic; $+$: phonon-dressed electronic; ad$^{\mathrm{s}}$: electron-following on the spin manifold of a magnet; qp: the quasiparticle-number sectors of a gapless target; $(2,0)$: pair breaking, and $(2,m)$ pair breaking with $m$ phonons. Columns: classes of coupling, with the target operators of \cref{sec:8.1}; ALP: axion-like particle; abs.: absorption. Entry: the one section that computes the rate; $\circ$: allowed at leading order, computed in no section; ---: no channel at leading order (\cref{sec:9.1-empty}). The density cell of single phonons is the coupling to charge through the Born charges; for a coupling to electrons alone, the acoustic peak of the electrons that follow the nuclei is computed in \cref{sec:10.2}. The absorption cells of excitons rest on measured optical data.}
  \label{tab:9-atlas}
  \begin{tabular}{@{}>{\raggedright\arraybackslash}p{0.31\linewidth}>{\centering\arraybackslash}p{0.068\linewidth}*{6}{>{\centering\arraybackslash}p{0.085\linewidth}}@{}}
    \toprule
    Excitation & Sector & Electron density & Nuclei & Spin, current & Vector abs. & ALP, scalar abs. & Inelastic, other \\
    & & $n_e$ & $n_\kappa$ & $\vect{s}$, $\vect{j}$, $\mathcal{J}$ & $\vect{j}_e$, $\vect{j}_\kappa$ & $\vect{s}_e$, $\mathcal{J}^{\alpha\alpha}_e$, $\mathcal{T}_e$ & rows of $\vect{q}$, $\Delta_\chi$ \\
    \midrule
    Electron--hole pairs, band to band (\ref{sec:3.2}, \ref{sec:4.2}) & 0 & \ref{sec:10.1} & --- & \ref{sec:10.4} & \ref{sec:10.7} & \ref{sec:10.7} & \ref{sec:10.9} \\
    Electron-gas continua and plasmons (\ref{sec:4.5}) & 0; qp & \ref{sec:10.3} & --- & --- & \ref{sec:10.7} & --- & \ref{sec:17.3} \\
    Excitons (\ref{sec:4.9}) & 0 & \ref{sec:10.5} & --- & \ref{sec:10.5} & \ref{sec:10.7} & \ref{sec:10.7} & --- \\
    Core and semicore electrons (\ref{sec:3.5}) & 0 & \ref{sec:10.6} & --- & $\circ$ & \ref{sec:10.7} & --- & --- \\
    \addlinespace[2pt]
    Single phonons (\ref{sec:5.3}, \ref{sec:5.5}) & ad & \ref{sec:11.2} & \ref{sec:11.1} & \ref{sec:11.6} & \ref{sec:11.8} & \ref{sec:11.8} & \ref{sec:17.3} \\
    Multiphonons (\ref{sec:5.6}) & ad & \ref{sec:11.3} & \ref{sec:11.3} & \ref{sec:11.6} & --- & --- & --- \\
    Free nuclear recoil (\ref{sec:1.3}, \ref{sec:5.6}) & ad & $\circ$ & \ref{sec:11.4} & \ref{sec:11.6} & --- & --- & \ref{sec:11.4} \\
    Quantum crystal: solid helium & ad & --- & \ref{sec:11.5} & --- & --- & --- & --- \\
    \addlinespace[2pt]
    Migdal effect: electrons with a nuclear recoil (\ref{sec:6.1}) & 0, $+$ & --- & \ref{sec:12.3} & $\circ$ & --- & \ref{sec:12.6} & --- \\
    Phonon-assisted electron recoil (\ref{sec:6.5}) & $+$ & \ref{sec:13.3} & --- & $\circ$ & \ref{sec:13.8} & \ref{sec:13.8} & --- \\
    \addlinespace[2pt]
    Magnons (\ref{sec:7.1}) & ad$^{\mathrm{s}}$ & --- & --- & \ref{sec:14.1} & \ref{sec:14.3} & \ref{sec:14.3} & \ref{sec:17.3} \\
    Two-magnon continua, magnon--phonon modes & ad$^{\mathrm{s}}$ & --- & --- & \ref{sec:14.2} & --- & --- & --- \\
    Chiral and nonreciprocal modes: crystals without inversion or time reversal & ad, ad$^{\mathrm{s}}$ & --- & --- & \ref{sec:14.5} & \ref{sec:14.5} & \ref{sec:14.5} & --- \\
    Precessing nuclear spins in a coherent field & ad & --- & --- & --- & --- & \ref{sec:14.6} & --- \\
    \addlinespace[2pt]
    Superconducting quasiparticles (\ref{sec:7.2}) & $(2,0)$, $(2,m)$ & \ref{sec:15.1} & --- & $\circ$ & \ref{sec:15.4} & \ref{sec:15.4} & --- \\
    Dirac and narrow-gap pairs (\ref{sec:7.3}) & qp; 0 & \ref{sec:16.1} & --- & $\circ$ & \ref{sec:16.1} & --- & --- \\
    Plasmons of sheets, wires and layered stacks & 0; qp & \ref{sec:16.2} & --- & --- & --- & --- & --- \\
    Graphene, nanotubes, quantum dots: confined electrons & 0; qp & \ref{sec:16.3} & --- & --- & \ref{sec:16.3} & --- & --- \\
    Frenkel excitons, localized on one molecule & 0 & \ref{sec:16.4} & --- & $\circ$ & --- & --- & --- \\
    Dopant transitions (\ref{sec:7.4}) & 0 & \ref{sec:16.5} & --- & $\circ$ & \ref{sec:16.5} & --- & --- \\
    Color centers and broken bonds (\ref{sec:7.4}) & ad & --- & \ref{sec:16.6} & --- & --- & --- & --- \\
    Two-level systems: tunneling atoms of a glass & ad & \ref{sec:16.7} & --- & --- & \ref{sec:16.7} & --- & --- \\
    Mobile defects: defectons and crowdions, kinetics only (\ref{sec:19.6}) & ad & --- & --- & --- & --- & --- & --- \\
    Acoustic modes of the whole body & ad & --- & --- & --- & \ref{sec:11.8} & \ref{sec:11.8} & --- \\
    \bottomrule
  \end{tabular}
\end{table}

\Cref{tab:9-atlas} is read in two directions. Along a row it lists the couplings that can create a given excitation, which is the question of a detector built to read out that excitation. Down a column it lists the excitations that a given coupling can create, which is the question of a model builder, and of \cref{chap:17}. Some entries stand for a short run of sections, and the table names the one that computes the rate. The density cell of the pair row, for instance, is the crystal form factor of independent electrons in \cref{sec:10.1}, and \cref{sec:10.2} screens it by the medium and writes it through the loss function. The nuclear cell of the Migdal row is the probability of \cref{sec:12.3}, written through the loss function of the crystal, which \cref{sec:12.2} derives in the atomic limit and \cref{sec:12.4} extends to phonon final states.

Coherent scattering off a whole random medium and the absorption of a dark photon by the vibrations of a glass, where no crystal momentum selects one optical phonon, are not single excitations and have no row; \cref{sec:16.7} treats them with the two-level systems. Processes second order in the coupling have no row either (\cref{sec:9.1-empty}).

The three entries of \cref{sec:17.3} in the last column are one formula, for exothermic or inelastic scattering into a collective mode, that serves phonons, magnons and plasmons alike, which is why the comparison chapter is their home. Couplings to the nuclear spin, and to baryon-minus-lepton number, which differs between isotopes through their neutron number, change from site to site. They reach the lattice through the random couplings of \cref{sec:11.6}. That section holds the spin entries of the lattice rows, from one phonon to a free recoil, where its self response at large $q$ becomes the impulse limit, and also the absorption of such bosons by phonons.

\subsection{Operators and sectors of the rows}\label{sec:9.1-sectors}

A row of \cref{tab:9-atlas} is a set of final states, and the decomposition theorem of \cref{sec:8.6} assigns every final state of a target with a gap to exactly one of three sectors. The sector of a row says which partial sum rule its rate shares and which other rows it can be added to. Two rows in different sectors add, apart from the interference cells of \cref{tab:8-sectors}; two cells in the same sector and the same column are pieces of one response, and their sections compute different parts of it, not terms of a sum.

The rows built on electronic excitations of an insulator sit in the zero-phonon electronic sector, $\mathcal{P}_0$. A band-to-band pair, an exciton, a core hole, a pair on a dopant and a Frenkel exciton all change the electronic state and leave the phonon number of the reference lattice unchanged at their zero-phonon line. The four electronic rows of the density column, pairs, plasmons, excitons and core electrons, are thus one response, $S^0_{ee}$, divided by the region of the $(\vect{q},\omega)$ plane and by the approximation that is accurate there. The exciton row is the part of the spectrum just above $E_{\mathrm{min}}(\vect{q})$ that the electron--hole attraction of \cref{sec:4.9} reshapes, the plasmon row the collective pole of \cref{sec:4.5}, and the core row the shells that open at large $\omega$ and $q$. The rate is the integral of $S^0_{ee}$ once, and the f-sum rule of \cref{sec:4.3} constrains the four cells together.

Phonons, multiphonons and a freely recoiling nucleus are all in the electron-following sector, $\mathcal{P}_{\mathrm{ad}}$: the electrons stay in their adiabatic ground state and the nuclei carry the energy. These three rows differ in the lattice regime alone of \cref{sec:8.5}, one phonon, several or a recoil, which grows with $q$ through the Debye--Waller exponent of \cref{sec:5.6}. The quantum crystal, the color centers created by a recoil, the two-level systems of a glass and the acoustic modes of a finite body are lattice states as well. Precessing nuclear spins belong to the same sector, although they are not a lattice state: the electrons stay in their ground state and the nuclear moments carry the energy. A color center is a different arrangement of the nuclei with the electrons again in their ground state; the electronic excitation that reads it out comes later and is not part of the dark-matter vertex. The electron-following sector is also how an electron coupling reaches the lattice. Through the density $\bar\rho_e(\vect{q};u)$ of electrons that follow the nuclei, \cref{eq:8.6-follow}, the electron-density column has entries in the single-phonon and multiphonon rows. In a polar crystal they enter through the electronic part of the Born effective charges, and at large momentum through the part of the electron cloud that a displaced atom carries along.

Two rows fill the phonon-dressed electronic sector, $\mathcal{P}_+$. The phonon-assisted row is the sideband of the electronic transitions, an electronic excitation with a net change of phonon number. The Migdal row straddles two sectors. The nuclear density reaches an electronic excitation only through the nonadiabatic boost of \cref{eq:8.6-boost}; its zero-phonon term lies in $\mathcal{P}_0$, and the terms in which the kick also creates phonons, or a free recoil, lie in $\mathcal{P}_+$. The interference of a Migdal amplitude with an electron recoil to the same final state is an off-diagonal block $S_{e\kappa}$ in these two sectors, nonzero for a coupling to both species (\cref{tab:8-sectors}). \Cref{sec:12.5} computes its zero-phonon block, and, with \cref{sec:17.4}, which treats the interference between channels, its phonon-dressed block.

The decomposition theorem needs a gap, and four kinds of row lack one in the sense of \cref{sec:8.6}. For a magnetic insulator the construction goes through with the low-energy spin manifold in place of the electronic ground state, so magnons and phonons share an electron-following sector, ad$^{\mathrm{s}}$ in the table. Superconducting pairs are counted by the number of Bogoliubov quasiparticles: pair breaking is the sector $(2,0)$, the analogue of $\mathcal{P}_0$, and pair breaking with a phonon, $(2, m \ge 1)$, the analogue of $\mathcal{P}_+$. For a metal, a Dirac semimetal and a gapless sheet such as graphene the quasiparticle-number sectors replace the three projectors, with the caveats of \cref{eq:8.6-bin}. A narrow-gap semiconductor and a semiconducting nanotube have a gap and use $\mathcal{P}_0$, which is why several rows carry two entries.

\subsection{Empty cells and selection rules}\label{sec:9.1-empty}

A cell is empty when the target operator of its column has no matrix element into the final states of its row at leading order. Leading order means first order in the coupling, one insertion of the interaction as in the master rate \cref{eq:8.2-rate}, and the lowest order of the expansion of \cref{sec:8.5} at which the process occurs. Operator content, the kinematics of absorption, the order in the coupling and the absence of a new response account for nearly all of the 79 empty cells.

\emph{Operator content.} The nuclear density cannot create an electronic excitation in the Born--Oppenheimer approximation. Its matrix element between the adiabatic ground state and an excited adiabatic state carries the electronic overlap $\langle\Psi_j(u)|\Psi_0(u)\rangle$, which vanishes at every nuclear configuration (\cref{sec:8.6}). The nuclear column is empty in every electronic row, and the nonadiabatic remainder, which does create pairs, is the Migdal row. In the same way the electron density cannot create a magnon. The density $n_e(\vect{q}) = \sum_j\ee^{\ii\vect{q}\cdot\vect{r}_j}$ acts on positions and commutes with the total electron spin, while a one-magnon state differs from the ordered ground state by one unit of the spin component along the order (\cref{sec:7.1}). Without spin--orbit coupling the matrix element vanishes, and the density column is empty in every magnetic row. The converse rule empties the spin column of the plasmon row. A spin density creates electron--hole pairs in the spin triplet, which carry no charge and produce no Coulomb field, and the plasmon is a pole that exists only because the induced charge feeds back through the long-range Coulomb interaction (\cref{sec:4.5}). The pair row has a spin entry, the plasmon row none.

\emph{Kinematics of absorption.} An absorbed boson brings $\omega = m_\phi$ at $q = m_\phi v \approx 10^{-3}m_\phi$. An acoustic phonon of that momentum has the energy $c_\nu m_\phi v$, with $c_\nu$ between $1.6\times10^{-5}$ and $3.1\times10^{-5}$ in silicon (\cref{sec:5.2}), some $10^{-8}$ of the boson mass. Since no acoustic phonon of an infinite crystal can absorb a boson, the single-phonon absorption entries are optical phonons at the zone center. A freely recoiling nucleus fares worse. Taking the momentum $m_\phi v$, it gains the energy $(m_\phi v)^2/2M$, a fraction $m_\phi v^2/2M$ of the boson mass. Exactly, a free particle at rest absorbs a boson of energy $E_\phi$ and momentum $\vect{p}_\phi$ only if $(E_\phi + M)^2 - p_\phi^2 = M^2$, that is $m_\phi^2 + 2E_\phi M = 0$, which no massive boson satisfies. A vector or axion-like coupling at $\vect{q} \to 0$ is linear in the nuclear velocities, so in a harmonic crystal it creates exactly one phonon. Two phonons need the nonlinear dependence of the dipole on the displacements, or anharmonicity, both higher orders of \cref{sec:8.5}, and the phase $\ee^{\ii\vect{p}_\phi\cdot\vect{u}}$ adds a correction of order $p_\phi^2\langle u^2\rangle$. The multiphonon and recoil rows are thus empty in the absorption columns. Two rows exist because these rules can be evaded. A body smaller than the wavelength of the boson has discrete acoustic modes, at MHz in a millimeter crystal, that a coherent field drives as a whole (\cref{sec:11.8}, which treats absorption on phonons and on the acoustic modes of a finite body). And a coupling that changes at random from site to site, such as the coupling to the nuclear spin of $^{29}$Si, does not conserve crystal momentum. A boson of negligible momentum can then create a phonon of any wavevector (\cref{sec:11.6}, which treats such random couplings).

\emph{Order in the coupling.} Two processes act on the target with two vertices: the absorption of two dark-matter particles at once, and dark Compton scattering, in which an absorbed boson leaves as a photon, $\phi e \to \gamma e$, with one dark and one electromagnetic vertex. They belong to no single cell and have no row (\cref{sec:10.8}, which builds their second-order transition operator).

\emph{No new response.} Inelastic or boosted dark matter changes where on the $(\vect{q},\omega)$ plane a target response is sampled, not the response. For most rows the rate of the last column follows from the elastic cell of the same operator with the kinematics of \cref{eq:8.1-omega-inel}, and the cell is left empty. The column has entries where the change of kinematics changes the physics. For pairs these are the endothermic threshold and the boosted and fermionic cases (\cref{sec:10.9}, which treats inelastic and relativistic dark matter). For nuclei it is fermionic absorption, where the outgoing fermion carries $q \simeq \mchi$ and the nucleus recoils with $\mchi^2/2M$ (\cref{sec:11.4}). For collective modes it is exothermic down-scattering at a fixed energy (\cref{sec:17.3}).

The row of mobile defects is empty by construction. Dark matter does not create defectons or crowdions at leading order; the row records that they release the energy of earlier excitations late, through the defect kinetics of \cref{sec:19.6}.

The eight cells marked $\circ$ are allowed at leading order, but no section computes their rate. In seven of them a spin coupling has a matrix element into the final states of the row. It can flip the spin of a core electron, break a Cooper pair or accompany a Migdal transition; it can create a Dirac pair or a phonon-assisted pair in the spin triplet; and it can flip the spin of a dopant electron or excite a molecule into its triplet state. The eighth is the electron density acting on a freely recoiling atom through the electrons that move with it (\cref{eq:8.6-follow}), whose rigid-atom density $f_\kappa(q)$ enters the multiphonon vertex $Z_\kappa - f_\kappa(q)$ of \cref{sec:11.3}. \Cref{sec:22.2}, which collects the open items of the book, lists the eight as one gap. A populated cell can also be empty for a particular material. Silicon has $Z^*_\kappa = 0$ (\cref{sec:5.5}): at long wavelength the electrons that follow each nucleus cancel its charge, so a dark photon couples to no single optical phonon of silicon, while in GaAs it couples through $Z^*_\kappa$, screened by $\varepsilon_\infty$. And a clean metal or superconductor, described as a free-electron gas, absorbs no dark photon below its interband transitions. The current at zero momentum is then the total momentum of the electrons divided by $m_e$, which the electron--electron interaction conserves, so absorption needs a third body, an impurity or a phonon; \cref{sec:15.4}, on absorption in superconductors and metals, computes both.

\subsection{Cells of a dark-matter model}\label{sec:9.1-models}

A heavy scalar mediator coupled to electrons enters only the density column: the electronic rows from pairs to core shells, the lattice through the electrons that follow the nuclei, the phonon-assisted row, and the superconductors and novel targets of \cref{chap:15,chap:16}. The same mediator coupled to nucleons enters only the nuclear column. It runs from single phonons through multiphonons to a free recoil, adds the Migdal row and creates color centers and broken bonds, and in silicon its lowest-mass coherent channel is the single acoustic phonon (\cref{sec:11.1}, which treats acoustic phonons). If its couplings to protons and neutrons differ, the isotopes of the crystal make it a random coupling, which reaches lighter masses incoherently (\cref{sec:11.6}). A spin-dependent coupling of the form $\mathcal{O}_4 = \vect{S}_\chi\cdot\vect{S}_e$ reaches the triplet pairs and excitons of a nonmagnetic crystal, and in a magnet the magnons and their continua. Coupled to nuclear spins instead, it reaches the lattice of an unpolarized crystal through the random couplings of \cref{sec:11.6}. Each of these models leaves most of \cref{tab:9-atlas} empty. The dark photon, which couples to every charge, does not.

\begin{workedexample}{The dark photon along its column}\label{we:9-darkphoton}
A dark photon couples to the electromagnetic current with strength $\zeta e$ (\cref{sec:1.4}), and \cref{tab:8-operators} reduces it to the current of all charges, $-\vect{j}_e + \sum_\kappa Z_\kappa\vect{j}_\kappa$. We follow it through \cref{tab:9-atlas}, first as the mediator of dark-matter scattering and then as the dark matter itself.

\emph{Scattering.} For a heavy dark photon the interaction is $-\Vol^{-1}\tilde V(q)\,n_Z(\vect{q})$, with $n_Z = -n_e + \sum_\kappa Z_\kappa n_\kappa$ the charge-number density (\cref{we:8-darkphoton}), so the dark photon enters the density and nuclear columns at once, with opposite signs. In the medium it mixes with the longitudinal photon, and the rate becomes an integral of the loss function of the total charge (\cref{sec:8.4}). That fixes its cells. All four electronic rows of the density column are open: the screened pair rate of \cref{sec:10.2}, which writes the electronic rate through the loss function, with its plasmon, exciton and core parts in \cref{sec:10.3,sec:10.5,sec:10.6}. In the lattice rows the electron and nuclear blocks combine into the charge carried by a vibration. An acoustic wave of long wavelength moves neutral cells and carries no charge, so the dark photon creates no acoustic phonon at small $q$, whereas an electron-only coupling does. An optical phonon of a polar crystal carries the Born charge, and this is the polar single-phonon rate of \cref{sec:11.2}; in silicon, with $Z^*_\kappa = 0$, the cell is empty. Above the zone the charge carried by a displaced atom is the nuclear charge minus the part of its cloud that follows, $Z_\kappa - f_\kappa(q)$, the charge vertex of the multiphonon response in \cref{sec:11.3}. Once $q$ exceeds the momenta of the core electrons the cloud is left behind, the vertex tends to $Z_\kappa$, and the dark photon makes the nucleus recoil freely with the weight $Z_\kappa^2$ (\cref{sec:11.4}). The same recoil creates color centers (\cref{sec:16.6}), and the same vertex, with the two electrons of helium, reaches the quantum crystal (\cref{sec:11.5}). Because the dark photon couples to both species, the Migdal row opens, together with its interference with electron recoil, the cross cells of \cref{tab:8-sectors}. The phonon-assisted row opens through its electron vertex (\cref{sec:13.3}, which builds the phonon-assisted structure factor), and so do the superconducting, Dirac, layered, confined, molecular, dopant and glass rows of the density column.

\emph{Absorption.} A dark photon of mass $m_\phi$ is absorbed at $q = m_\phi v$ and $\omega = m_\phi$, and with the current row it measures the loss function of the target at that point. The vector column of \cref{tab:9-atlas} lists where this is nonzero. On electrons it is the optical response at $\omega = m_\phi$ (\cref{sec:10.7}), up through the L and K edges of the core shells; near the plasma frequency the dark photon converts into the plasmon. In silicon between the indirect gap of \SI{1.17}{eV} and the direct gap of \SI{3.4}{eV} no vertical transition exists (\cref{sec:3.3}), and absorption needs a phonon, the phonon-assisted cell of \cref{sec:13.8}. On the lattice the absorber is the zone-center longitudinal optical phonon of a polar crystal, at $m_\phi = \omega_{\mathrm{L}} = \SI{36.35}{meV}$ in GaAs (\cref{we:5-lst}), and nothing in silicon; a finite piezoelectric body adds its acoustic modes, which the field inside it drives through the polarization of a strain (\cref{sec:11.8}). In a magnet the dark photon reaches the spins through its magnetic field, $\vect{B}' = \ii\vect{p}_\phi\times\hat{\vect{e}}\,A'$, with $\hat{\vect{e}}$ the polarization and $A'$ the field amplitude. That field is smaller than the electric one by the factor $v$ (\cref{sec:14.3}, which treats absorption into magnons), and single-molecule magnets absorb in the same way (\cref{sec:14.5}). A superconductor needs an impurity or a phonon to absorb (\cref{sec:15.4}), and the remaining entries, Dirac materials, aligned nanotubes, dopants and two-level systems, each measure their own loss function (\cref{sec:16.1,sec:16.3,sec:16.5,sec:16.7}).

\emph{The last column.} Inelastic dark matter that scatters through a dark photon uses the same charge density with the kinematics of \cref{eq:8.1-omega-inel}. Its endothermic and exothermic pair rates are in \cref{sec:10.9}, and its exothermic down-scattering into plasmons and optical phonons in \cref{sec:17.3}.

Altogether the dark photon populates all 14 cells of the density column, all six of the nuclear column, all 14 of the vector column and three of the last: 37 of the 57 channels of the atlas.
\end{workedexample}
